%% Question bank chapter file - generated by bank_tools.py build (do not edit while the bank is still being built:
%% it is regenerated from the per-test data in Question Bank\data). Every question is one bankq block:
%%   \begin{bankq}[id=..., source={...}, paper=.., qno=.., type=..., subtopic={...}, answer={...}, answer_src=key|solved,
%%                 working={...}, difficulty=1-10, diff_parts={K S I C F}, diff_note={...},
%%                 group=<passage id or empty>, also_in={duplicate copies in other tests}]
%% difficulty = min(10, K+S+I+C+F-1): K concepts 1-3, S steps 1-4, I insight 0-2, C computation 0-2, F format 0-1
%%   ...question LaTeX, options, figures (\bfig / \bfigi -> img/)...
%%   \end{bankq}
%% Passages shared by several questions are bankpassage blocks; their questions carry group=<passage id>.
\documentclass[10pt]{article}
\usepackage[a4paper,margin=1.8cm]{geometry}
\usepackage[T1]{fontenc}
\usepackage{amsmath,amssymb,mathtools}
\usepackage{graphicx}
\usepackage[export]{adjustbox}
\usepackage{xcolor,needspace,keyval}
\usepackage{enumitem}
\usepackage{booktabs,array,tabularx}
\usepackage[hidelinks]{hyperref}
\setlength{\parindent}{0pt}
\setlength{\parskip}{3pt}
\makeatletter
\define@key{bq}{id}{\def\bq@id{#1}}\define@key{bq}{source}{\def\bq@src{#1}}\define@key{bq}{paper}{\def\bq@paper{#1}}
\define@key{bq}{qno}{\def\bq@qno{#1}}\define@key{bq}{type}{\def\bq@type{#1}}\define@key{bq}{subtopic}{\def\bq@sub{#1}}
\define@key{bq}{answer}{\def\bq@ans{#1}}\define@key{bq}{answer_src}{\def\bq@asrc{#1}}\define@key{bq}{working}{\def\bq@work{#1}}
\define@key{bq}{group}{\def\bq@group{#1}}\define@key{bq}{also_in}{\def\bq@also{#1}}
\define@key{bq}{difficulty}{\def\bq@diff{#1}}\define@key{bq}{diff_parts}{\def\bq@dparts{#1}}\define@key{bq}{diff_note}{\def\bq@dnote{#1}}
\newenvironment{bankq}[1][]{\def\bq@work{}\def\bq@also{}\def\bq@group{}\def\bq@diff{?}\def\bq@dparts{}\def\bq@dnote{}\setkeys{bq}{#1}%
  \par\needspace{4\baselineskip}\medskip\noindent{\small\textbf{\bq@id}\hfill\color{gray}[\bq@src, Paper \bq@paper, Q\bq@qno]
  \bq@type\ $\cdot$ \bq@sub\ $\cdot$ \textbf{D\bq@diff}}\par\nopagebreak\relax}
  {\par{\small\color{teal}Answer: \textbf{\bq@ans}\ifx\bq@work\@empty\else\ (solved: \bq@work)\fi
  \ifx\bq@also\@empty\else\par\color{gray}Also in: \bq@also\fi}\par}
\define@key{bp}{id}{\def\bp@id{#1}}\define@key{bp}{source}{\def\bp@src{#1}}\define@key{bp}{questions}{\def\bp@qs{#1}}
\newenvironment{bankpassage}[1][]{\setkeys{bp}{#1}\par\needspace{4\baselineskip}\bigskip\noindent
  {\small\textbf{Passage \bp@id}\hfill\color{gray}[\bp@src; questions \bp@qs]}\par\nopagebreak\relax\itshape}{\par}
\makeatother
\newcommand{\bfig}[2]{\par\smallskip{\centering\includegraphics[width=#1pt,max width=0.62\linewidth,max height=0.32\textheight]{img/#2}\par}\smallskip}
\newcommand{\bfigi}[2]{\includegraphics[width=#1pt,max width=\linewidth,max height=0.32\textheight,valign=c]{img/#2}}
\newcommand{\bmissing}[1]{\textit{[figure missing in the source conversion: \detokenize{#1}]}}
\newenvironment{qoptions}{\begin{itemize}[leftmargin=1.8em,itemsep=1pt,topsep=2pt,label={}]}{\end{itemize}}
\newcommand{\opt}[2]{\item[(#1)] #2}
\providecommand{\fwstart}{}\providecommand{\fwend}{}\providecommand{\qfullsplit}{}
\newcolumntype{L}[1]{>{\raggedright\arraybackslash}p{#1}}
\newcolumntype{C}[1]{>{\centering\arraybackslash}p{#1}}
\newcolumntype{R}[1]{>{\raggedleft\arraybackslash}p{#1}}
\providecommand{\sectionhdr}[1]{}\providecommand{\subsectionhdr}[1]{}


\begin{document}
{\Large\bfseries Thermodynamics --- Question Bank}\par\smallskip
{\small Chemistry chapter C04. Questions from all converted test papers, grouped by source series. Answers: printed key where it matches the paper, otherwise solved (answer\_src=solved, with working). Exact repeats across papers appear once, with also\_in listing the other copies.}\par\medskip\hrule

\section*{Score-I papers}
\begin{bankq}[id=SI12-P1-CHEM-S2-Q02, source={Score-I - Test 12 (09 Mar 2022)}, paper=1, qno=2, type=numerical, subtopic={Laws}, answer={137.76}, answer_src=solved, working={$V_1=24.6$ L at 1 atm, 300 K; step I doubles $T$ to 600 K; adiabatic II back to 300 K: $V_X=V_1(600/300)^{1/(\gamma-1)}=24.6\times2^{5/2}=24.6\times5.6\approx137.8$ L}, difficulty=6, diff_parts={K2 S3 I0 C1 F1}, diff_note={Heating at const V (T doubles), adiabatic T V^(gamma-1) relation with gamma=1.4, V1=nRT/P}]
Initially one mole of ideal gas $\left(C_v=\dfrac52R\right)$ at 1.0 atm and 300 K is put through the following cycle:\\
Step-I: Heating to twice its initial pressure at constant volume. ($R=0.082\left(\dfrac{\text{litre-atm}}{\text{K-mole}}\right)$, $\sqrt2=1.4$)\\
Step-II: Adiabatic expansion to its initial temperature.\\
Step-III: Isothermal compression back to 1.00 atm.\\
What is the volume at state $X$ (in litre)?
\bfig{71.1}{SI12_p1_chem_s2_q02_1.png}
\end{bankq}
\begin{bankq}[id=SI12-P2-CHEM-S2-Q05, source={Score-I - Test 12 (09 Mar 2022)}, paper=2, qno=5, type=numerical, subtopic={Gibbs energy}, answer={7.75}, answer_src=solved, working={$\Delta C_p=0$, so for $A+2B\to3C$: $\Delta G=20-300(0.015)=15.5$ kJ; half the reaction: 7.75 kJ/mol}, difficulty=5, diff_parts={K1 S3 I1 C0 F1}, diff_note={Delta Cp=0 so dH, dS constant; dG=20000-300*15 for full reaction, halve it}]
Molar heat capacities at constant pressure for $A$, $B$ and $C$ are 3, 1.5 and 2 J/K mol. The enthalpy of reaction and entropy of reaction, $A+2B\to3C$ are 20 kJ/mol and 15 J/K-mole at 300 K. Calculate $\Delta G$ (in kJ/mol) for the reaction.\\
$\dfrac12A+B\to\dfrac32C$
\end{bankq}
\end{document}
